\part{Market Failures}

\chapter{Externalities}
\label{chap:externalities}

\section{Learning objectives}

After completing this chapter, you should be able to:
\begin{itemize}
  \item define positive and negative externalities;
  \item distinguish private from social costs and benefits;
  \item identify the inefficient market quantity; and
  \item compare taxes, subsidies, property rights, regulation, and permits.
\end{itemize}

\section{External effects}

An \keyterm{externality} occurs when an action affects people not directly involved in the market transaction and this effect is not reflected in the market price. A factory's emissions may harm neighbours; a vaccination may reduce infection risk for others. These are respectively negative and positive externalities.

With a negative production externality, marginal social cost exceeds marginal private cost:
\[
MSC=MPC+MEC,
\]
where $MEC$ is marginal external cost. Competitive firms choose output using $MPC$, so the market produces too much relative to the efficient quantity, where marginal social benefit equals marginal social cost.

With a positive externality, marginal social benefit exceeds marginal private benefit. The market produces too little because buyers and sellers do not capture all benefits.

\section{Policy responses}

A \keyterm{Pigouvian tax} equal to marginal external cost makes private decision-makers face the relevant social cost. A subsidy can encourage activities with positive external benefits. Direct regulation can prohibit harmful activities or set standards. Tradable permits can limit total pollution while allowing firms with low abatement costs to reduce more.

The \keyterm{Coase theorem} highlights that, if property rights are clear and bargaining is costless, parties can negotiate an efficient outcome regardless of who initially holds the right. In practice, transaction costs, many affected parties, unequal bargaining power, and information problems often limit bargaining.

\section*{Chapter Review}

\subsection*{Key concepts}
\begin{itemize}
  \item \textbf{Externality}: an uncompensated effect on a third party.
  \item \textbf{Negative externality}: $MSC>MPC$.
  \item \textbf{Positive externality}: $MSB>MPB$.
  \item \textbf{Pigouvian tax}: tax designed to internalise external cost.
  \item \textbf{Coase theorem}: bargaining can achieve efficiency under restrictive conditions.
\end{itemize}

\section*{Exercises}

\subsection*{Conceptual questions}
\begin{enumerate}[label=\arabic*.]
  \item Define an externality and give one positive and one negative example.
  \item Why does a negative production externality create excessive market output?
  \item State the main conditions required for Coasean bargaining.
\end{enumerate}

\subsection*{Graphical exercises}
\begin{enumerate}[resume,label=\arabic*.]
  \item Draw demand, MPC, and MSC for a negative production externality. Mark market output, efficient output, and deadweight loss.
  \item Show how a Pigouvian tax changes the firm's private marginal-cost curve.
\end{enumerate}

\subsection*{Analytical exercises}
\begin{enumerate}[resume,label=\arabic*.]
  \item Demand is $P=100-Q$, private marginal cost is $20$, and external marginal cost is $20$. Find market and efficient quantities and the corrective tax.
  \item Explain when a tradable-permit system may be preferable to a uniform technology standard.
\end{enumerate}

\subsection*{Exam-style multiple-choice questions}
\begin{enumerate}[resume,label=\arabic*.]
  \item A negative externality exists when: (a) private cost exceeds social cost; (b) social cost exceeds private cost; (c) all firms make losses; or (d) price is zero.
  \item A corrective tax on pollution is intended to: (a) make pollution free; (b) align private and social costs; (c) increase the efficient quantity; or (d) eliminate all industry.
\end{enumerate}

\subsection*{PPE discussion question}

\noindent Should climate policy prioritise carbon taxes, regulation, or public investment? Discuss the efficiency, equity, and political-feasibility dimensions of the choice.
